\documentclass{article}
\title{How amplitude changes a pendulum's period}
\author{Terso example manuscript}
\date{}
\begin{document}
\maketitle
\begin{abstract}
A pendulum released from ten degrees takes about 0.19 percent longer to
complete a cycle than the small-angle model predicts.
\end{abstract}

\section{Model}
Let $\theta_0$ denote the release angle in radians. For an ideal pendulum of
length $L$ in uniform gravity $g$, the small-angle period is
$T_0=2\pi\sqrt{L/g}$. The exact period is proportional to the complete
elliptic integral with modulus $\sin(\theta_0/2)$. Expanding that integral
using \cite{dlmf19} gives Eq.~(\ref{eq:period}).

% !TEX root = ../main.tex
\section{Results}
\begin{equation}
  \frac{T}{T_0} = 1 + \frac{\theta_0^2}{16} + O(\theta_0^4).
  \label{eq:period}
\end{equation}
At ten degrees, the leading correction is $(\pi/18)^2/16$, or about
0.19 percent. Table~\ref{tab:angles} shows how this estimate changes with
the release angle.

\begin{table}[h]
\centering
\begin{tabular}{rc}
Release angle (degrees) & Leading increase in period (percent) \\
5 & 0.048 \\
10 & 0.190 \\
15 & 0.428
\end{tabular}
\caption{Fractional increase over the small-angle period from the quadratic
term in Eq.~(\ref{eq:period}). Values are rounded.}
\label{tab:angles}
\end{table}


\bibliographystyle{plain}
\bibliography{references}
\end{document}
